\chapter{Chapter 4: Serving, Scaling, and Optimizing LLMs}

\section*{Context}

\textbf{Problem statement:} Engineer low-latency, high-throughput LLM serving with predictable TTFT/TPOT, bounded cost/query, and graceful degradation under traffic spikes.

\subsection*{Real-World Examples}

\begin{itemize}
  \item A coding assistant launch hits 12x expected traffic; continuous batching and KV caching reduce TTFT from 2.1s to 900ms without additional GPUs.
  \item A support assistant outage reveals retry storms from provider 429s; queue isolation and deterministic degrade modes restore stability within minutes.
\end{itemize}

\bigskip\hrule\bigskip

\section*{Step 1: Understand the Problem and Establish Design Scope}

\subsection*{Simulated Dialogue}

\textbf{Interviewer:} Design the serving and scaling infrastructure for a multi-tenant LLM platform.\\
\textbf{Candidate:} I need to understand the traffic shape, latency SLOs, model mix, and cost targets before proposing architecture.

\textbf{Candidate:} What is the request profile?\\
\textbf{Interviewer:} 20M requests/day, mix of short completions (avg 200 output tokens) and long generations (avg 800 output tokens). Peak is 15x average.

\textbf{Candidate:} What latency SLOs must we hit?\\
\textbf{Interviewer:} Interactive routes: TTFT $<$ 800ms, TPOT $<$ 40ms. Batch routes: p95 completion $<$ 30s.

\textbf{Candidate:} What is the model landscape?\\
\textbf{Interviewer:} Three tiers: 7B quantized (fast), 34B (balanced), 70B (complex). Plus managed API fallback.

\subsection*{Functional Requirements}

\begin{itemize}
  \item Serve multiple model tiers (7B, 34B, 70B) with automatic routing by request complexity and SLO.
  \item Support streaming responses with token-level delivery.
  \item Handle managed API failover when self-hosted capacity is exhausted.
  \item Enforce per-tenant and per-route token budgets and rate limits.
\end{itemize}

\subsection*{Non-Functional Requirements}

\begin{itemize}
  \item \textbf{Latency:} TTFT p95 $<$ 800ms (interactive), TPOT p95 $<$ 40ms/token.
  \item \textbf{Throughput:} sustain 3,500 QPS peak with $<$ 5\% error rate.
  \item \textbf{Cost:} self-hosted unit cost $<$ 60\% of equivalent managed API cost at steady state.
  \item \textbf{Availability:} 99.9\% uptime with graceful degradation under overload.
  \item \textbf{GPU utilization:} target $>$ 70\% average utilization across serving fleet.
\end{itemize}

\subsection*{Back-of-the-Envelope Estimation}

\subsubsection*{Traffic Profile}

\begin{itemize}
  \item Requests/day: 20M
  \item Average QPS: $20M / 86{,}400 \approx 231$
  \item Peak QPS: $\sim$3,470 (15x average)
  \item Token mix: 60\% short (200 output tokens), 40\% long (800 output tokens)
  \item Average output tokens: $0.6 \times 200 + 0.4 \times 800 = 440$ tokens/request
  \item Output tokens/day: $20M \times 440 = 8.8B$ tokens
\end{itemize}

\subsubsection*{GPU Capacity Planning (Self-Hosted)}

\begin{itemize}
  \item 7B quantized (INT4): $\sim$18GB VRAM, $\sim$120 tokens/sec/GPU
  \item 34B quantized (INT8): $\sim$40GB VRAM, $\sim$45 tokens/sec/GPU
  \item 70B (FP16, tensor parallel across 2xA100): $\sim$160GB VRAM, $\sim$25 tokens/sec/GPU
\end{itemize}

\textbf{Peak throughput demand (tokens/sec):}

\begin{itemize}
  \item Peak QPS $\times$ avg output tokens: $3{,}470 \times 440 \approx 1.53M$ tokens/sec
  \item Route split at peak: 50\% fast (7B), 35\% balanced (34B), 15\% complex (70B)
  \item 7B demand: $765K$ tok/sec $/\, 120 \approx 6{,}375$ GPU-equivalents $\rightarrow$ with continuous batching (4--8x efficiency): \textbf{$\sim$800--1,600 GPUs}
  \item 34B demand: $\sim$535K tok/sec $/\, 45 \approx 11{,}900$ GPU-equiv $\rightarrow$ with batching: \textbf{$\sim$1,500--3,000 GPUs}
  \item 70B demand: $\sim$230K tok/sec $/\, 25 \approx 9{,}200$ GPU-equiv $\rightarrow$ with batching: \textbf{$\sim$1,150--2,300 GPUs}
  \item Total fleet: \textbf{$\sim$3,500--7,000 GPUs} (order-of-magnitude for peak; autoscaling reduces steady-state to $\sim$30\% of this)
\end{itemize}

\subsubsection*{KV Cache Memory}

\begin{itemize}
  \item 70B model, 128K context window, 80 layers, 128 heads, 128 head-dim:
  \begin{itemize}
    \item Per-token KV: $2 \times 80 \times 128 \times 128 \times 2$ bytes (FP16) $\approx 5.2$MB/token
    \item 4K context: $\sim$20.5GB KV cache per request
    \item Concurrent batch of 16 requests: $\sim$328GB $\rightarrow$ requires careful memory management
  \end{itemize}
\end{itemize}

\subsubsection*{Cost Comparison}

\begin{itemize}
  \item Managed API (GPT-4 class): $\sim$\$30/1M input tokens, $\sim$\$60/1M output tokens
  \item Self-hosted 70B (amortized GPU + ops): $\sim$\$8--12/1M output tokens at scale
  \item Breakeven: self-hosted wins when sustained volume $>$ $\sim$2M output tokens/day per GPU
\end{itemize}

\bigskip\hrule\bigskip

\section*{Step 2: Propose High-Level Design and Get Buy-In}

\subsection*{Core Architecture Components}

\begin{itemize}
  \item \textbf{Load Balancer + API Gateway:} TLS termination, tenant auth, rate limiting, request classification.
  \item \textbf{Request Classifier:} estimates complexity/length to select model tier and queue priority.
  \item \textbf{Model Router:} routes to self-hosted tier or managed API fallback based on capacity and SLO.
  \item \textbf{Serving Runtime (per tier):} continuous batching scheduler, KV cache manager, streaming output handler.
  \item \textbf{Autoscaler:} GPU fleet autoscaling based on queue depth, utilization, and latency signals.
  \item \textbf{Cache Hierarchy:} prefix/KV cache (VRAM), semantic response cache (Redis), session prompt cache.
  \item \textbf{Managed API Gateway:} rate-limited client for external provider failover.
  \item \textbf{Observability:} TTFT/TPOT histograms, GPU utilization, queue depth, cache hit rate, error rate by tier.
\end{itemize}

\subsection*{Pipeline A: Model Deployment Pipeline}

\begin{enumerate}
  \item Model artifact (weights + tokenizer + config) versioned in model registry.
  \item Quantization and optimization applied (GPTQ/AWQ for INT4/INT8, or FP16 for large tier).
  \item Serving config generated: max batch size, max sequence length, KV cache budget per GPU.
  \item Canary deployment to shadow traffic (5\% of production load).
  \item Quality gate: compare output quality, latency, and error rate against baseline.
  \item Gradual traffic shift (5\% $\rightarrow$ 25\% $\rightarrow$ 50\% $\rightarrow$ 100\%) with automatic rollback triggers.
\end{enumerate}

\subsection*{Pipeline B: Request Serving Pipeline}

\begin{enumerate}
  \item Request arrives at gateway; auth, rate limit, and quota checks applied.
  \item Request classifier estimates token count and complexity $\rightarrow$ assigns model tier and priority.
  \item Model router checks self-hosted capacity (queue depth $<$ threshold); falls back to managed API if overloaded.
  \item Serving runtime adds request to continuous batch scheduler.
  \item Scheduler groups requests by sequence length, manages KV cache allocation, begins prefill.
  \item Tokens streamed to client as generated; TTFT and TPOT tracked per request.
  \item On completion: response logged, telemetry emitted, cache updated if eligible.
\end{enumerate}

\textit{[Insert Architecture Diagram: Load balancer $\rightarrow$ request classifier $\rightarrow$ model router branching to self-hosted tiers (7B/34B/70B each with continuous batching runtime and KV cache) and managed API fallback. Autoscaler watches queue depth and GPU utilization. Cache hierarchy at multiple levels. Observability collecting TTFT/TPOT/utilization.]}

\subsection*{Mermaid Architecture Diagram}

\begin{verbatim}
flowchart LR
  client[Client] --> lb[LoadBalancer]
  lb --> gateway[APIGateway]
  gateway --> classifier[RequestClassifier]
  classifier --> router[ModelRouter]

  router --> fast[FastTier7B]
  router --> balanced[BalancedTier34B]
  router --> complex[ComplexTier70B]
  router --> managed[ManagedAPIFallback]

  fast --> batchSched1[ContinuousBatchScheduler]
  balanced --> batchSched2[ContinuousBatchScheduler]
  complex --> batchSched3[ContinuousBatchScheduler]

  batchSched1 --> kvCache1[KVCacheManager]
  batchSched2 --> kvCache2[KVCacheManager]
  batchSched3 --> kvCache3[KVCacheManager]

  kvCache1 --> stream[StreamingOutput]
  kvCache2 --> stream
  kvCache3 --> stream
  managed --> stream

  stream --> client

  autoscaler[Autoscaler] --> fast
  autoscaler --> balanced
  autoscaler --> complex

  semanticCache[SemanticResponseCache] --> router
  observability[Observability] --> autoscaler
  stream --> observability
\end{verbatim}

\bigskip\hrule\bigskip

\section*{Step 3: Design Deep Dive}

Focus areas for this chapter:

\begin{enumerate}
  \item \textbf{Continuous Batching and Scheduling}
  \item \textbf{KV Cache Management}
  \item \textbf{Autoscaling and Overload Protection}
\end{enumerate}

\subsection*{1) Continuous Batching and Scheduling}

Traditional static batching wastes GPU cycles waiting for the longest sequence in a batch to complete. \textbf{Continuous batching} (also called iteration-level or in-flight batching) admits new requests and retires completed ones at every decode step.

\textbf{How it works:}

\begin{itemize}
  \item Scheduler maintains a running batch of active requests.
  \item At each decode iteration, completed requests are evicted and new requests are admitted.
  \item Prefill (prompt processing) and decode (token generation) can be interleaved or separated.
\end{itemize}

\subsubsection*{Batching Strategy Trade-off Table}

\small
\begin{tabular}{|p{0.16\linewidth}|p{0.15\linewidth}|p{0.15\linewidth}|p{0.14\linewidth}|p{0.24\linewidth}|}
\hline
\textbf{Strategy} & \textbf{GPU Utilization} & \textbf{TTFT Impact} & \textbf{Complexity} & \textbf{Best For} \\
\hline
Static batching & 30--50\% & High (waits for batch fill) & Low & Dev/testing \\
\hline
Continuous batching & 60--80\% & Low (admit immediately) & Medium & Production default \\
\hline
Chunked prefill & 70--85\% & Very low (prefill interleaved) & High & High-throughput serving \\
\hline
Disaggregated prefill/decode & 80--90\% & Lowest & Very high & Hyperscale platforms \\
\hline
\end{tabular}
\normalsize

\textbf{Scheduling policies:}

\begin{itemize}
  \item \textbf{FCFS (first-come first-served):} simple, fair, but head-of-line blocking from long prompts.
  \item \textbf{Shortest-job-first (SJF):} lower average latency but starvation risk for long requests.
  \item \textbf{Priority queues:} route-level isolation, SLO-based priority, with starvation timeout.
  \item \textbf{Recommended:} priority queues with SLO-based preemption and starvation guard (max wait 5s).
\end{itemize}

\subsection*{2) KV Cache Management}

The KV cache stores computed key/value attention states for previously processed tokens. It is the primary memory bottleneck in LLM serving.

\textbf{Memory anatomy for 70B model, single request, 4K context:}

\begin{itemize}
  \item Per-layer KV: $2 \text{ (K+V)} \times \text{head\_dim} \times \text{num\_heads} \times 2$ bytes (FP16)
  \item 80 layers, 128 heads, 128 head-dim: $\sim$5.2MB/token
  \item 4K context: $\sim$20.5GB per concurrent request
  \item Batch of 32: $\sim$656GB $\rightarrow$ exceeds single-node capacity
\end{itemize}

\subsubsection*{KV Cache Optimization Trade-off Table}

\small
\begin{tabular}{|p{0.18\linewidth}|p{0.14\linewidth}|p{0.17\linewidth}|p{0.14\linewidth}|p{0.23\linewidth}|}
\hline
\textbf{Technique} & \textbf{Memory Savings} & \textbf{Quality Impact} & \textbf{Complexity} & \textbf{Best For} \\
\hline
No optimization (FP16 full) & Baseline & None & None & Small models, low concurrency \\
\hline
KV cache quantization (INT8) & $\sim$50\% & Minimal ($<$0.5\% quality loss) & Low & Default for large models \\
\hline
PagedAttention (vLLM) & Variable (reduces fragmentation) & None & Medium & Production serving engines \\
\hline
Sliding window attention & Proportional to window size & Degrades on long-range dependencies & Medium & Long-context with locality \\
\hline
Multi-query / grouped-query attention & 4--8x per head & Built into model architecture & None (model-level) & Newer model architectures \\
\hline
\end{tabular}
\normalsize

\textbf{PagedAttention (vLLM):}

\begin{itemize}
  \item Manages KV cache as virtual memory pages (blocks of token slots).
  \item Eliminates fragmentation from pre-allocated contiguous buffers.
  \item Enables memory sharing for beam search and parallel sampling.
  \item Production default in vLLM, TensorRT-LLM, and similar serving engines.
\end{itemize}

\textbf{Prefix caching:}

\begin{itemize}
  \item Cache KV states for common system prompts and instruction prefixes.
  \item If system prompt is 2K tokens and repeats across 90\% of requests, prefix cache saves $\sim$10.4GB recomputation per request (70B model).
  \item TTFT reduction: skip prefill for cached prefix $\rightarrow$ TTFT drops by 200--400ms.
  \item Invalidation: version-tag prefixes; invalidate on prompt template change.
\end{itemize}

\subsection*{3) Autoscaling and Overload Protection}

\subsubsection*{Autoscaling Signals}

\small
\begin{tabular}{|p{0.22\linewidth}|p{0.17\linewidth}|p{0.12\linewidth}|p{0.35\linewidth}|}
\hline
\textbf{Signal} & \textbf{Scaling Direction} & \textbf{Lag} & \textbf{Notes} \\
\hline
Queue depth (requests waiting) & Scale up & Seconds & Most responsive signal \\
\hline
GPU utilization & Scale up/down & Minutes & Smoothed average; avoid thrashing \\
\hline
TTFT p95 & Scale up & Minutes & Directly tied to user experience \\
\hline
Request error rate (timeouts, OOM) & Scale up (emergency) & Immediate & Circuit-breaker integration \\
\hline
Cost budget remaining & Scale down / cap & Hours & Prevents runaway spend \\
\hline
\end{tabular}
\normalsize

\textbf{Overload protection mechanisms:}

\begin{itemize}
  \item \textbf{Admission control:} reject or defer requests when queue exceeds threshold (e.g., 2x batch capacity).
  \item \textbf{Load shedding:} drop low-priority requests first; serve cached/degraded responses for others.
  \item \textbf{Circuit breaker to managed API:} when self-hosted error rate $>$ 5\%, overflow to managed provider.
  \item \textbf{Token budget enforcement:} hard-reject requests exceeding max input/output token limits.
  \item \textbf{Graceful degradation:} reduce max output tokens, skip reranking, serve cache-only summaries.
\end{itemize}

\subsection*{Production Bottlenecks and Mitigations}

\begin{itemize}
  \item \textbf{GPU OOM during long-context requests}
  \begin{itemize}
    \item Symptom: CUDA OOM errors, request failures on long prompts.
    \item Mitigation: dynamic max sequence length per batch, prompt truncation policy, KV cache quantization.
  \end{itemize}

  \item \textbf{Queue saturation and head-of-line blocking}
  \begin{itemize}
    \item Symptom: TTFT spikes for all requests when one long prompt dominates.
    \item Mitigation: priority queues, chunked prefill, sequence-length-based bucketing.
  \end{itemize}

  \item \textbf{Provider 429 storms during failover}
  \begin{itemize}
    \item Symptom: managed API rate limits hit during self-hosted outage; cascading failure.
    \item Mitigation: token bucket rate limiting on managed API client, jittered retries, request prioritization.
  \end{itemize}

  \item \textbf{KV cache fragmentation}
  \begin{itemize}
    \item Symptom: effective batch size drops despite available VRAM.
    \item Mitigation: PagedAttention, periodic cache compaction, right-sizing page block size.
  \end{itemize}

  \item \textbf{Autoscaler thrashing}
  \begin{itemize}
    \item Symptom: rapid scale up/down cycles during variable traffic.
    \item Mitigation: cooldown periods, smoothed signals, min fleet size floor.
  \end{itemize}

  \item \textbf{Model version mismatch during canary}
  \begin{itemize}
    \item Symptom: inconsistent outputs during partial rollout.
    \item Mitigation: sticky routing by session during canary window, clear version labeling in responses.
  \end{itemize}
\end{itemize}

\subsection*{Advanced Techniques}

\begin{itemize}
  \item \textbf{Speculative decoding:} small draft model proposes tokens, large model verifies in parallel. 2--3x speedup when acceptance rate $>$ 70\%.
  \item \textbf{Continuous batching with chunked prefill} for TTFT stability under mixed workloads.
  \item \textbf{LoRA adapter serving:} load domain-specific adapters on shared base model, hot-swap per request. Adds $\sim$5\% overhead per adapter.
  \item \textbf{Tensor parallelism} for models exceeding single-GPU VRAM (70B+ class).
  \item \textbf{Quantized KV cache (INT8)} for 50\% memory savings with negligible quality loss.
\end{itemize}

\bigskip\hrule\bigskip

\section*{Step 4: Wrap Up}

\subsection*{Final Design Summary}

The design treats serving as latency engineering: route-tiered inference with continuous batching, hierarchical KV cache management, and deterministic overload controls that prioritize user experience under stress.

\subsection*{Wrap-Around Concerns}

\begin{itemize}
  \item \textbf{Release engineering:} canary model deployments with quality gates and automatic rollback.
  \item \textbf{Governance:} model version tracking, GPU allocation audits, cost attribution per tenant/route.
  \item \textbf{Cost controls:} autoscaler cost caps, managed API budget limits, cache hit-rate targets ($>$ 30\% for semantic cache).
  \item \textbf{Drift controls:} latency and quality dashboards sliced by model version, route, and traffic cohort.
\end{itemize}

\bigskip\hrule\bigskip

\section*{Step 5: Interviewer Curveballs}

\subsection*{Interviewer Follow-ups \& Frequently Asked Questions (FAQ)}

\textbf{Q: TTFT jumped from 700ms to 2.2s overnight. Walk me through your debugging steps.}

\textbf{The Ideal Answer:} Check queue depth first (did traffic spike or did capacity drop?). Then check cache hit rates (did a deployment invalidate prefix cache?). Then check for routing shifts (are more requests hitting the slow tier?).

\textbf{Deep Dive \& Trade-offs:}
\begin{itemize}
  \item Queue depth rise = capacity problem or traffic spike.
  \item Cache hit-rate drop = deployment-induced invalidation or traffic pattern change.
  \item Tier shift = request classifier regression or model router misconfiguration.
\end{itemize}

\textbf{The Interviewer's Rubric:} Checking structured incident triage for LLM-specific serving issues.

\textbf{Q: How do you decide when to add a new model tier vs scaling an existing one?}

\textbf{The Ideal Answer:} When latency distribution is bimodal (simple requests penalized by sharing queue with complex ones), add a smaller tier. When distribution is unimodal, scale horizontally.

\textbf{Deep Dive \& Trade-offs:}
\begin{itemize}
  \item New tier: better SLO isolation but more operational complexity (deployment, monitoring, routing).
  \item Horizontal scaling: simpler ops but wastes GPU on over-provisioned model for simple requests.
\end{itemize}

\textbf{The Interviewer's Rubric:} Checking fleet architecture reasoning beyond ``just add GPUs.''

\textbf{Q: What is the practical limit of continuous batching efficiency?}

\textbf{The Ideal Answer:} Bounded by KV cache memory. Maximum concurrent requests $\approx$ \texttt{available\_VRAM / (per\_request\_KV\_size)}. At 80GB VRAM and 2.5GB KV/request (34B, 2K context), max batch $\approx$ 32.

\textbf{Deep Dive \& Trade-offs:}
\begin{itemize}
  \item KV quantization (INT8) doubles effective batch size.
  \item PagedAttention reduces fragmentation overhead by 10--20\%.
  \item Trade-off: larger batches improve throughput but increase per-token latency (TPOT).
\end{itemize}

\textbf{The Interviewer's Rubric:} Checking whether you understand the memory-throughput-latency triangle.

\textbf{Q: When is speculative decoding worth the complexity?}

\textbf{The Ideal Answer:} When the draft model acceptance rate exceeds $\sim$70\% and TPOT reduction matters more than throughput (interactive use cases, not batch).

\textbf{Deep Dive \& Trade-offs:}
\begin{itemize}
  \item Draft model adds VRAM overhead ($\sim$2--4GB for 1B draft model).
  \item Low acceptance rate ($<$50\%) can actually increase latency due to wasted compute.
  \item Best for: long generations where user is streaming and perceives token-by-token speed.
\end{itemize}

\textbf{The Interviewer's Rubric:} Checking nuanced understanding of inference optimization trade-offs.

\textbf{Q: How do you handle a tenant that sends 10K-token prompts when your system is optimized for 2K average?}

\textbf{The Ideal Answer:} Route to dedicated long-context queue with separate batch scheduling, enforce maximum context limit per tier, and charge premium pricing for long-context usage.

\textbf{Deep Dive \& Trade-offs:}
\begin{itemize}
  \item Long prompts consume 5x KV cache, blocking other requests in shared batch.
  \item Separate queue prevents head-of-line blocking.
  \item Consider prompt summarization as pre-processing step if full context is not needed.
\end{itemize}

\textbf{The Interviewer's Rubric:} Checking multi-tenancy fairness and resource isolation thinking.

\textbf{Q: What GPU utilization target should you set, and why not 100\%?}

\textbf{The Ideal Answer:} Target 70--80\% average. Headroom absorbs traffic bursts without queueing. At $>$90\%, any spike causes immediate latency degradation.

\textbf{Deep Dive \& Trade-offs:}
\begin{itemize}
  \item 100\% utilization = zero burst absorption = queueing on any traffic fluctuation.
  \item 60\% utilization = wasted capacity and poor economics.
  \item Autoscaler should maintain 70--80\% band with scale-up triggered at 85\%.
\end{itemize}

\textbf{The Interviewer's Rubric:} Checking capacity planning discipline.

\textbf{Q: How do you cost-compare self-hosted vs managed API for a specific route?}

\textbf{The Ideal Answer:} Calculate all-in self-hosted cost (GPU amortization + networking + ops labor + electricity) per 1M output tokens, and compare against managed API per-token pricing at projected steady-state volume.

\textbf{Deep Dive \& Trade-offs:}
\begin{itemize}
  \item Self-hosted: high fixed cost, low marginal cost. Wins at sustained high volume.
  \item Managed: zero fixed cost, high marginal cost. Wins for bursty/low-volume routes.
  \item Breakeven typically at $\sim$2--5M output tokens/day per GPU depending on model size.
\end{itemize}

\textbf{The Interviewer's Rubric:} Checking total cost of ownership reasoning.

\textbf{Q: How would you serve LoRA adapters for 50 different tenants on shared infrastructure?}

\textbf{The Ideal Answer:} Load base model once per GPU. Hot-swap LoRA adapters per request using adapter registry. Use adapter batching to group requests with same adapter.

\textbf{Deep Dive \& Trade-offs:}
\begin{itemize}
  \item Each LoRA adapter: $\sim$50--200MB (rank 16--64, applied to attention layers).
  \item Adapter switch overhead: $\sim$2--5ms per request (negligible).
  \item Challenge: cache pollution if adapter diversity is high and batch sizes are small.
  \item Alternative: tenant-dedicated model instances for high-volume tenants.
\end{itemize}

\textbf{The Interviewer's Rubric:} Checking parameter-efficient serving architecture knowledge.

\section*{End of Chapter}

\begin{itemize}
  \item \textbf{Key takeaway:} LLM serving is memory-bound, not compute-bound. KV cache management, batching strategy, and overload protection determine whether your system hits latency SLOs at scale.
  \item \textbf{What to practice:} Size a GPU fleet for a specific traffic profile with back-of-the-envelope math including VRAM, throughput, and cost.
  \item \textbf{Interview drill:} Given a TTFT regression, walk through the full debug path from queue depth to cache hit rate to model routing to GPU utilization.
\end{itemize}
